% Generated draft: requires mathematical and source review.
\documentclass[11pt]{article}
\usepackage[margin=25mm]{geometry}
\usepackage{fontspec}
\setmainfont{Latin Modern Roman}
\usepackage{amsmath,amssymb,mathtools,unicode-math}
\setmathfont{Latin Modern Math}
\usepackage{enumitem,needspace,xurl,hyperref}
\setlength{\parindent}{0pt}
\setlength{\parskip}{5pt}
\setlength{\emergencystretch}{3em}
\newcommand{\N}{\mathbb{N}}
\newcommand{\Z}{\mathbb{Z}}
\newcommand{\Q}{\mathbb{Q}}
\newcommand{\R}{\mathbb{R}}
\newcommand{\C}{\mathbb{C}}
\newcommand{\F}{\mathbb{F}}
\newcommand{\A}{\mathbb{A}}
\newcommand{\PP}{\mathbb{P}}
\newcommand{\E}{\mathbb{E}}
\newcommand{\eps}{\varepsilon}
\newcommand{\dd}{\,\mathrm{d}}
\begin{document}
\section*{UnsolvedMath ID: 30001254}
\textbf{Source problem code:} catalogue entry
\subsection{Definitions and mathematical statement}
Let \(K\) be a field complete with respect to a discrete valuation, and let \(F\) be a one-variable function field over \(K\). The usual finitely generated-extension reading is that \(F/K\) is finitely generated and \(\operatorname{trdeg}_K F=1\); the supplied record does not specify any additional constant-field condition.

Let \(\Omega_F\) denote the set of equivalence classes of nontrivial rank-one discrete valuations of \(F\). For \(v\in\Omega_F\), write \(F_v\) for the completion of \(F\) at \(v\). The quantifier ranges over all such valuations: the supplied statement imposes no restriction on their restrictions to \(K\).

Let \(G\) be a connected linear algebraic group over \(F\), and let \(Y\) be a homogeneous space under \(G\), meaning an \(F\)-variety with a \(G\)-action that is transitive after extension to an algebraic closure of \(F\). For a field extension \(L/F\), let \(Y(L)\) denote its set of \(L\)-rational points.

\textbf{Question.} For every such \(K,F,G,Y\), does the implication
\[
\bigl(\forall v\in\Omega_F,\quad Y(F_v)\ne\varnothing\bigr)
\quad\Longrightarrow\quad
Y(F)\ne\varnothing
\]
hold?

\textbf{Unresolved source qualification.} The accompanying supplied assessment describes the question as removing an \(F\)-rationality hypothesis from a result in the report, but it does not identify the object subject to that hypothesis or reproduce the result's full assumptions. The displayed question preserves the supplied problem statement, which imposes no \(F\)-rationality hypothesis. Its precise relationship to the report's qualified result requires human source review. No restrictions on characteristic, residue field, or stabilizers are stated in the supplied formulation.
\subsection{Short English statement}
The question asks whether a homogeneous space of a connected linear algebraic group over a one-variable function field over a complete discretely valued field has a rational point whenever it has a point over every rank-one discrete valuation completion. The supplied background mentions removing an \(F\)-rationality hypothesis without specifying that hypothesis precisely. This is a statement-only draft requiring human mathematical review.
\subsection{Sources}
\begin{itemize}
\item Local UnsolvedMath registry, permanent ID \(30001254\), source problem code \texttt{catalogue entry}, title \emph{Local--Global Principles for Homogeneous Spaces}, in the \emph{Oberwolfach Reports Open Problems} collection. The supplied record is the local registry evidence for this draft; repository registry locations are \texttt{data/unsolvedmath/problems.json} and \texttt{catalog/problems.json}.
\item Supplied source citation: \emph{Quadratic Forms and Linear Algebraic Groups} (2009), Oberwolfach Reports, DOI 10.4171/owr/2009/25, pp.~1375--1426. Original Oberwolfach report: \url{https://doi.org/10.4171/owr/2009/25}. The registry gives proposed year 2009 and names no individual proposer. Its background describes Oberwolfach Reports as documenting mathematical workshops and their open-problem discussions.
\item Supplied literature reference: Linh, \emph{arithmetics of homogeneous spaces over p-adic function fields}, \url{https://londmathsoc.onlinelibrary.wiley.com/doi/full/10.1112/jlms.12842}. No fuller bibliographic attribution is supplied.
\item Supplied literature assessment dated exactly \texttt{2026-08-21}: the record labels the problem open and says that the report leaves open removal of an \(F\)-rationality hypothesis, while later results address special group and stabilizer classes rather than the full statement. This is supplied dated curation, not an independently verified assessment. The supplied background expressly calls for independent verification of time-sensitive status and cited claims.
\item The registry labels the statement \texttt{corrected\_verified} and reports that its mathematical content was checked against the official Oberwolfach report, with notation, context, and question form normalized without changing the source-level claim. These are supplied verification claims; no source visit or independent mathematical verification is asserted here.
\end{itemize}
\end{document}
